\section{Compatibility, stability contrasts and interface checks}\label{sec:compatibility}
\subsection{Why local static equality is not global causal attainment}
If an incoming row has law $\mu_s$, any $m$-label randomized assignment has within-row distortion at least $q_m(\mu_s)$. In~\eqref{eq:three-defect}, equality in this row bound and vanishing pooling require, respectively, rowwise optimal quantization and agreement of all incoming centroids assigned to a common outgoing label. These are \emph{local} conditions. They do not assert the existence of one recursion satisfying them through every preceding information cut.

\begin{example}[Two explicit obstructions]\label{ex:compatibility}
Two equally likely rows $\delta_0$ and $\delta_1$ each have $q_1=0$. Merging them into one outgoing label has pooling distortion $1/4$, despite rowwise optimality. For an obstruction at an earlier cut, acquire two independent fair signs $B_1,B_2$, and only afterwards receive an independent fair index $I\in\{1,2\}$ requesting $B_I$. A two-label checkpoint after $I$ can answer exactly. A two-label state before $I$ cannot: its two possible readout pairs are a two-center quantizer for the four equally likely vertices of $\{-1,1\}^2$. The optimal total squared distortion is one, by checking adjacent two--two partitions versus one--three partitions; division by the independent index gives risk at least $1/2$. A randomized encoder cannot improve the pointwise minimum over its two response pairs. Thus a small final answer alphabet is not a global recursive memory certificate.
\end{example}

A full-dimensional terminal law is also insufficient for a uniform row condition. In the neutral circle-rotation experiment, a full-history predictive vector can be uniform on the unit circle while a one-label retained barycenter is zero. Its subsequent rotation row is then $\delta_0$. This is why the lower theorem quantifies over all legal retained barycenters, not merely a typical full-history law. The example is outside Definition~\ref{def:class}, and is used as an assumption witness, not a realization of its density theorem.

\subsection{Integrated curvature and same-report stability are different}
For comparison with the retained Bayesian theorem, call (C) the existence of block-summable upper curvature coefficients for the perspective potential, and call (P) block-summable pathwise Lipschitz coefficients for every same-report normalized update in a legal metric. Neither is silently substituted for the new acquired-overlap condition.

\begin{example}[(C) fails while (P) holds]\label{ex:Cfails}
A binary sign has prior mean $u$, a noisy sign report has density $(1+\eta yu)/2$, and a positive symmetric hidden transition multiplies the posterior mean by $\lambda$. Thus
$F(y,u)=\lambda(u+\eta y)/(1+\eta yu)$, $y\in\{-1,1\}$. Direct differentiation gives
\[
 \tfrac12\Phi''(u)=
 \lambda^2(1-\eta^2)^2\frac{1+3\eta^2u^2}{(1-\eta^2u^2)^3}.
\]
At $\lambda=\eta=u=.99$ this is greater than $24$. The homogeneous uniform (C) coefficient therefore cannot have summable products. In the metric $2|\operatorname{arctanh}u-\operatorname{arctanh}v|$, the report adds a constant to the arctanh coordinate, while $u\mapsto\lambda u$ has derivative ratio
$\lambda(1-u^2)/(1-\lambda^2u^2)\le\lambda<1$.
Thus (P) holds on the positive post-transition interval. The example separates stability routes; its finite report alphabet does not satisfy the full-dimensional row-density hypothesis.
\end{example}

\begin{example}[(P) fails while (C) holds]\label{ex:Pfails}
With probability $\omega\in(0,1)$ a marked quantum branch does nothing and reports its mark. Otherwise it performs the covariant rank-one instrument
$D_qP_\psi\rho P_\psi\,d\nu(\psi)$, where $\nu$ is uniform projective measure. Use the centered Hilbert--Schmidt predictive chart. The pure posterior has constant squared norm $1-1/D_q$, so the integrated potential is
\[
 \Phi(\rho-I/D_q)=\omega\norm{\rho-I/D_q}_{\HS}^2
                         +(1-\omega)(1-1/D_q).
\]
Its curvature coefficient is exactly $\omega<1$. On the identity branch every nondegenerate metric has same-report Lipschitz coefficient at least one, at every step. Hence no uniform pathwise (P) product contracts. All branch masses are real and positive; the closed perspective covers possible zero likelihood at an individual pure state. This finite-rank example is a contrast for the two sufficient routes, not an assertion of the interior/density hypotheses of Theorem~\ref{thm:main}.
\end{example}

\subsection{Clarifications for the retained passive companion}
The passive Gaussian convention in the retained Bayesian article is report first, positive transition second; its constants include the determinant lower of the affine likelihood-rank matrix and $\alpha\le1/q$. Its positive-cone contraction constants depend on $\lambda,D_q$, and its bounded-gap extension starts in a declared interior domain before the first improving step. In a projective-orbit small-ball estimate, a ball centered off the orbit either misses it or is contained, on that orbit, in a ball of twice the radius centered at an orbit point. Finite-stratum equal allocation is a sufficient fixed-weight upper, not a weight-optimal allocation theorem.

In that companion's \emph{passive} early-bit extension, the core remains independent of the bit even after conditioning on all its allowed reports and its fixed action word. A completely forgotten fair sign has squared loss one for responses in $[-1,1]$. The shrinkage response $(1-a)b$ in its two-lane construction has risk $2a-a^2$; with $\varepsilon=a/2$ this is at most $4\varepsilon$. If $Q=M-2n-1\ge2$, its core budget $M_0=1+\lfloor(M-3)/2\rfloor$ obeys
$M_0-n=\lfloor Q/2\rfloor\ge Q/3$; the remaining bounded values use its stated coarse construction. These are clarifications of the retained interface. They do not transfer the pre-core bit lower to arbitrary feedback actions. Theorem~\ref{thm:joint} instead proves the controlled post-core version with a different, explicit state count.

All finite-time simulation comparisons include the committed answer, target coins and protected audit. Loss bounds retain their physical-dimension normalization. These conventions rule out a free stored response, a latent external clock, or a simulator that reads the audit before generating its preceding report.
